% !PRINTED 323
\weberpage{323}
%<<CONT
angles change their sign as often as the point $x$ passes through a side
of the coordinate triangle.

With the help of this non-homogeneous relation, unity can be
represented linearly and homogeneously by $x_1, x_2, x_3$, and thereby
every non-homogeneous function can be transformed into a homogeneous
one.

According to these determinations, the coordinates $x_1, x_2, x_3$
are at first not independent variables, but are bound to the relation
mentioned. But we now extend their meaning so that we can regard
them as independent variables, if we stipulate that we shall use them
only in homogeneous forms or equations;
and that not $x_1, x_2, x_3$ themselves,
but only the ratios $x_1 : x_2 : x_3$ shall determine the position of the point $x$.
Then we can disregard entirely the relation mentioned
between the $x_1, x_2, x_3$, and regard the $x_1, x_2, x_3$ as independent variables, for which only the single combination of values $x_1 = 0, x_2 = 0, x_3 = 0$ is excluded. The
values $h x_1, h x_2, h x_3$ then determine, for every $h$ different from zero, one and the same point $x$. Strictly speaking, the
coordinates $x_1, x_2, x_3$ of the point $x$ are then not the distances
from the coordinate axes described above themselves, but any quantities whatever that are proportional to these distances.

Every linear substitution

\weq{1}{\begin{aligned}
x_1 &= \alpha_1^{(1)} y_1 + \alpha_1^{(2)} y_2 + \alpha_1^{(3)} y_3 \\
x_2 &= \alpha_2^{(1)} y_1 + \alpha_2^{(2)} y_2 + \alpha_2^{(3)} y_3 \\
x_3 &= \alpha_3^{(1)} y_1 + \alpha_3^{(2)} y_2 + \alpha_3^{(3)} y_3,
\end{aligned}}

\noindent whose determinant

\weq{2}{r = \Sigma \pm \alpha_1^{(1)} \alpha_2^{(2)} \alpha_3^{(3)}}

\noindent is different from zero, can be interpreted geometrically in two
ways: first as a \gesperrt{mapping}, in that we regard $x_1, x_2, x_3$
and $y_1, y_2, y_3$ as coordinates of two points $x, y$ in the same
coordinate system, so that to every point $y$ there corresponds a point $x$
and conversely; and then also as a \gesperrt{coordinate transformation}, if we regard $x_1, x_2, x_3$ and $y_1, y_2, y_3$ as coordinates of one and the same point, referred to two different
coordinate systems. For the considerations that come next
we shall hold to the latter interpretation.
